\subsection{EXAMPLES OF RINGS}

I. Zero ring. Let A be a ring. For 0 = 1 in A, it is necessary and sufficient that A consist of a single element. The condition is obviously sufficient. On the other hand, if 0 = 1, then, for all \(x \in A\) , x = x.1 = x.0 = 0. Such a ring is called a zero ring.

\begin{enumerate}
\def\labelenumi{\Roman{enumi}.}
\setcounter{enumi}{1}
\tightlist
\item
  Ring of rational integers. With the addition defined in §2, no. 5, and the multiplication defined in §2, no. 6, Z is a commutative ring. The notation 0, 1, -x is in accordance with the notation introduced earlier.
\end{enumerate}

*III. Ring of real-valued functions. Let I be an interval in the set R of real numbers and let A be the set of continuous functions defined on I with real values. The sum \(f + g\) and product f.g of two functions f and g are defined by

\[
(f + g) (t) = f (t) + g (t), \quad (f g) (t) = f (t) g (t) \quad (t \in \mathbf {I}).
\]

A commutative ring is obtained whose unit element is the constant 1.*

*IV. Convolution pseudo-ring. Let E be the set of real-valued continuous functions on R, which are zero outside a bounded interval. The sum of two functions is defined as in III, but the product is now defined by

\[
(f g) (t) = \int_ {- \infty} ^ {\infty} f (s) g (t - s) d s
\]

(``convolution product''). Thus a commutative pseudo-ring is obtained which is not a ring (cf.~Integration, VIII, §4).*

V. Opposite ring of a ring A. Let A be a ring. The set A with the same addition as A and the multiplication \((x,y)\mapsto yx\) is often denoted by \(A^{0}\) . It is a ring (called the opposite ring of A) with the same zero and same unit as A and which coincides with A if and only if A is commutative.

\begin{enumerate}
\def\labelenumi{\Roman{enumi}.}
\setcounter{enumi}{5}
\tightlist
\item
  Endomorphism ring of a commutative group. Let G be a commutative group written additively. Let E denote the set of endomorphisms of G. Given f and g in E, the mappings \(f + g\) and fg of G into G are defined by
\end{enumerate}

\[
(f + g) (x) = f (x) + g (x), \quad (f g) (x) = f (g (x)) \quad (x \in \mathbf {G}).
\]

By § 1, no. 5, Proposition 5, \(f + g\) is an endomorphism of G and so obviously also is \(fg = f \circ g\) . By § 4, no. 8, E is a (commutative) group under addition. Multiplication is obviously associative and has identity element \(\operatorname{Id}_{\mathbf{G}}\) . Also for \(f, g\) and \(h\) in E, we write \(\phi = f \cdot (g + h)\) ; for all \(x \in \mathbf{G}\) ,

\[
\phi (x) = f ((g + h) (x)) = f (g (x) + h (x)) = f (g (x)) + f (h (x))
\]

for \(f\) is an endomorphism of \(\mathbf{G}\) ; hence \(\phi = fg + fh\) and clearly

\[
(g + h) f = g f + h f.
\]

Therefore E is a (not in general commutative) ring called the endomorphism ring of G.

\begin{enumerate}
\def\labelenumi{\Roman{enumi}.}
\setcounter{enumi}{6}
\tightlist
\item
  Pseudo-ring of zero square. A pseudo-ring A is said to be of zero square if xy = 0 for all \(x, y \in A\) . Let G be a commutative group. If the set G is given the addition of the group G and multiplication \((x, y) \mapsto 0\) , a pseudo-ring of zero square is obtained. It is only a ring if \(G = \{0\}\) , in which case it is the zero ring.
\end{enumerate}

